\subsection{Packing independent sets}\label{finite:packing-subsection}

The forbidden parameters now give lower bounds on independent sets in
exterior balls.  Throughout this subsection, $G$ has no $C_{k+1}$,
$\omega(G)\le t$, and $\alpha(G)\le k$, where $k\ge5$ and $3\le t\le k$.
We also use the expansion property
\[
 |N_G[I]|\ge k|I|+1
 \qquad\text{for every nonempty independent set }I,
\]
which was established in Lemma~\ref{red:expansion}.
In particular, $\delta(G)\ge k$.
Let $S\subseteq X\subseteq V(G)$, and let $M$ be a symmetric matrix of
valid forbidden outside-path parameters on $X$, with empty diagonal,
as defined above.  Write $n_X=|X|$.  The balls $B_r(i;X)$ are those of
Section~\ref{sec:paths}.

\begin{lemma}[Bounds for exterior balls]\label{finite:ball-bounds}
For $i\in X$, put
\[
 b_0=k+1-n_X+|\{j\in X:0\in M_{ij}\}|.
\]
Assume $b_0>0$, and define
\[
 u_0=1+\mathbf 1_{\{b_0\ge t\}}.
\]
For $r=1,2$, recursively set
\begin{align}
 A_r(i)&=\{j\in X:\{1,\ldots,r\}\subseteq M_{ij}\},
 &T_r&=|A_r(i)|,\label{finite:excluded-neighbors}\\
 b_r&=\max\{b_{r-1},\,ku_{r-1}+1-n_X+T_r\},\label{finite:ball-size}\\
 \widehat u_r&=\max\bigl\{u_{r-1},
       1+\mathbf 1_{\{b_r>t\}}
        +\mathbf 1_{\{b_r>\max(k,2t)\}}\bigr\}.
       \label{finite:ordinary-weight}
\end{align}
Set $u_r=\widehat u_r$, except that $u_r=2$ when
\begin{equation}\label{finite:exception-condition}
 \widehat u_r=1,\qquad b_{r-1}=b_r=t,\qquad
 k+1-n_X+T_r\ge t.
\end{equation}
Then, for $r=0,1,2$,
\[
 |B_r(i;X)|\ge b_r,\qquad \alpha(B_r(i;X))\ge u_r.
\]
\end{lemma}

\begin{proof}
The empty diagonal means that the vertices counted by
$|\{j:0\in M_{ij}\}|$ are all distinct from $i$.  They are nonneighbors
of $i$.  Thus at most
$n_X-1-|\{j:0\in M_{ij}\}|$ neighbors of $i$ lie in $X$, proving the
bound on $|B_0(i;X)|$.  Positivity gives an independent vertex.
Furthermore, a $t$-clique in $B_0(i;X)$ together with $i$ would be a
$(t+1)$-clique.  Therefore $b_0\ge t$ forces an independent pair, which
proves the asserted value of $u_0$.

Suppose the bounds have been proved at radius $r-1$.  Each vertex
$j\in A_r(i)$ is distinct from $i$, because $M_{ii}$ is empty, and has
no neighbor in $B_{r-1}(i;X)$.  Indeed, such a neighbor would give an
outside $i$--$j$ path with between one and $r$ internal vertices, by
the no-neighbor assertion of Lemma~\ref{path:separation}.
Choose an independent set $I\subseteq B_{r-1}(i;X)$ of order
$u_{r-1}$.  Its closed neighborhood outside $X$ is contained in
$B_r(i;X)$, whereas its closed neighborhood inside $X$ avoids the
$T_r$ vertices of $A_r(i)$.  Expansion gives
\[
 |B_r(i;X)|\ge ku_{r-1}+1-(n_X-T_r).
\]
The balls are nested, so \eqref{finite:ball-size} follows.
Nesting also preserves the independence bound $u_{r-1}$.
If $b_r>t$, the ball is not a clique; if $b_r>\max(k,2t)$,
Lemma~\ref{size:test} supplies an independent triple.
These are precisely the additional bounds in
\eqref{finite:ordinary-weight}, with the stated strict inequalities.

It remains to justify \eqref{finite:exception-condition}.
If $B_r(i;X)$ were a clique $C$, nesting and
$b_{r-1}=b_r=t$ would force
\[
 B_{r-1}(i;X)=B_r(i;X)=C,\qquad |C|=t.
\]
Every $v\in C$ has at most $t-1$ neighbors outside $X$, since
$B_r(i;X)$ contains all outside neighbors of the preceding ball.
The vertices of $A_r(i)$ are also unavailable as neighbors.  Hence
\[
 k+1-t
 \ \le\ |N_G(v)\cap X|
 \ \le\ n_X-T_r
 \ \le\ k+1-t.
\]
If the last inequality were strict, there would already be a
contradiction.  Otherwise equality forces $v$ to be adjacent to every
vertex of $X\setminus A_r(i)$, including $i$.
Thus $C\cup\{i\}$ is a $(t+1)$-clique, a contradiction.
The ball therefore contains an independent pair.
\end{proof}

For each $i\in X$, include the \emph{singleton option} $\{i\}$, label its
radius by $-1$, and give it weight one.  If $b_0>0$, also include the
options $B_r(i;X)$, for $r=0,1,2$, with weights $u_r$ from
Lemma~\ref{finite:ball-bounds}.  If $b_0\le0$, include no ball option at
$i$: the available bound does not certify that a ball is nonempty.
This omission makes no assertion about the actual ball.

\begin{lemma}[Packing test]\label{finite:packing}
Let a family of these options have distinct base vertices.
For options based at $i,j$ with radius labels $r,s$, require
\begin{equation}\label{finite:compatibility}
 \begin{cases}
 0\in M_{ij},&r=s=-1,\\
 \{1,\ldots,r+s+2\}\subseteq M_{ij},&\text{otherwise}.
 \end{cases}
\end{equation}
If every pair satisfies this condition, the sum of the option weights
is at most $k$.
\end{lemma}

\begin{proof}
For $r,s\ge0$, the balls are disjoint and anticomplete by
Lemma~\ref{path:separation}.  For a singleton $\{i\}$ and a ball
$B_s(j;X)$, disjointness is automatic.  An edge between them would give
an outside $i$--$j$ path with between one and $s+1$ internal vertices,
which is prohibited by \eqref{finite:compatibility} with $r=-1$.
Two singleton options are nonadjacent precisely when their joining edge
is absent, as certified by $0\in M_{ij}$.
Thus all represented sets are disjoint and pairwise anticomplete.
Independent subsets of their certified weights combine into an
independent set in $G$.  Its order cannot exceed $k$.
\end{proof}

\paragraph{Removing redundant radii.}
Retain every singleton and every available radius-zero option.
For $r=1,2$, retain the option at radius $r$ only when
$u_r>u_{r-1}$; still compute both bounds at every radius.
This does not change whether a family of weight greater than $k$
exists.  Indeed, an omitted positive radius $r$ has the same weight as
the most recent retained \emph{nonnegative} radius $q<r$ at that base.
For every other radius $s\ge-1$, replacing $r$ by $q$ weakens the
positive-interval condition in \eqref{finite:compatibility}.
Since $q+s+2\ge1$, this replacement never invokes the exceptional
singleton--singleton condition.
Replacing every omitted option in a compatible family therefore
preserves its weight and compatibility.  Distinct bases ensure that
no replacement duplicates another member of the family.

\paragraph{Searching for a packing.}
Make a finite weighted graph whose vertices are the retained options,
with edges given by \eqref{finite:compatibility} between distinct bases.
The packing test asks for a clique of total weight greater than $k$.
The following inclusion-and-exclusion search answers that question
exactly.  For a candidate list $V$ and a nonnegative threshold $K$,
take its first member $x$.  If its weight $w(x)$ exceeds $K$, return
success.  Otherwise let $Z$ be the remaining candidates adjacent to
$x$.  The branch containing $x$ succeeds exactly when $Z$ contains a
clique of weight greater than $K-w(x)$; the branch excluding $x$ is
the same question on $V\setminus\{x\}$.
The inclusion branch can be discarded when
$\sum_{z\in Z}w(z)\le K-w(x)$.
The empty list fails.  Induction on $|V|$ proves the recursion correct:
every clique either contains $x$ or avoids it, and the discarded branch
cannot reach the required weight.  All recursive thresholds are
nonnegative, since the one-option success was tested first.
Any ordering of the options gives the same decision.

\subsection{Strengthening the forbidden parameters}
\label{finite:strengthening-subsection}

We apply the packing test on the original set $S$ and, temporarily,
on a set with one additional vertex.  On either set, a required edge
whose parameter zero is forbidden is an immediate contradiction.
The other test is a compatible family of weight greater than $k$.

\begin{lemma}[Eliminating a path hypothesis]\label{finite:strengthening}
Let $J$ be a required-edge graph on $S$, and let $M$ be a matrix of
valid forbidden parameters relative to $S$.
Fix distinct $i,j\in S$ and $d\in\{0,1\}$.
Form a trial required-edge graph as follows:
\begin{enumerate}
\item If $d=0$, retain $X=S$ and require the edge $ij$ in addition to $J$.
\item If $d=1$, add a fresh vertex $z$, put $X=S\cup\{z\}$, and require
the edges $iz,zj$ in addition to $J$.
\end{enumerate}
On the trial set, retain the old entries of $M$ on pairs in $S$,
start with empty entries involving $z$ when it is present, and add all
forbidden parameters from Lemma~\ref{finite:required-path-rule}.
If these trial data give either an immediate required-edge
contradiction or a packing contradiction, then $d$ is forbidden
between $i$ and $j$ relative to $S$.
\end{lemma}

\begin{proof}
Assume an outside $i$--$j$ path with parameter $d$ exists relative to
$S$.  For $d=0$ it is the edge $ij$, so the trial required-edge graph
is realized on $S$.
For $d=1$, choose its actual internal vertex $z\notin S$.
The trial required-edge graph is then realized on $S\cup\{z\}$.
There may be other edges involving $z$; the required graph does not
assert their absence.

In both cases, every old forbidden parameter remains valid on $X$
by Lemma~\ref{finite:inheritance}: restricting the allowed outside
vertices cannot create a previously forbidden path.
The additional required-path constraints are valid by
Lemma~\ref{finite:required-path-rule}.
Either stipulated contradiction is therefore impossible in this
realization.  This excludes the assumed path.
In the case $d=1$, the argument applies to any actual choice of its
internal vertex, so the conclusion is a prohibition relative to the
original set $S$, not merely to one enlarged trial set.
\end{proof}

The rules just proved give the following finite procedure for an
optimal path-system pattern $P$.
\begin{enumerate}
\item Construct its required graph $J$ on $S$ and its initial matrix
$M$ by Lemma~\ref{finite:extension-rule}.  If the packing test gives a
contradiction, return $1$.
\item Add all required-path constraints of
Lemma~\ref{finite:required-path-rule}.  If there is a required-edge or
packing contradiction, return $2$.
\item Using the current matrix unchanged, test every unordered pair
$i,j\in S$ and every $d\in\{0,1\}$ not already in $M_{ij}$.
Skip $d=0$ when $ij$ is already required by $J$.
For each remaining choice, perform the trial of
Lemma~\ref{finite:strengthening}, and collect the choices whose trial
gives a contradiction.
\item If none were collected, return $0$.
Otherwise add every collected parameter to both symmetric entries of
$M$.  Return $2$ if the updated matrix gives a required-edge or packing
contradiction; if it does not, repeat Step~3.
\end{enumerate}
All tests in Step~3 use the same matrix from the beginning of that
step.  The required graph $J$ on $S$ never changes.  Consequently its
required-path constraints need only be added once in Step~2.

\begin{proposition}[Meaning of a successful finite check]
\label{finite:procedure-soundness}
The procedure terminates.  An output of $1$ or $2$ proves that the
pattern $P$ cannot be realized by an optimal path system in a graph
with the hypotheses above.  An output of $0$ asserts only that these
tests have not produced a contradiction.
\end{proposition}

\begin{proof}
Suppose such a realization exists.  Its initial forbidden parameters
are valid by Lemma~\ref{finite:extension-rule}, and the constraints
added in Step~2 are valid by Lemma~\ref{finite:required-path-rule}.
Inductively assume the current matrix contains valid prohibitions.
Each new entry collected in Step~3 follows from that same matrix by
Lemma~\ref{finite:strengthening}.  Every collected prohibition
therefore holds in the realization, and they may all be added
simultaneously.  No new entry has been assumed in proving another
entry of its batch.
This proves the invariant at every subsequent step.
An immediate required-edge contradiction or a violation of
Lemma~\ref{finite:packing} is thus impossible, proving the assertion
about outputs $1$ and $2$.

Every pass through Step~4 that does not terminate adds a previously
absent parameter from the finite set of
$2\binom{|S|}{2}$ pair-and-parameter choices.
The procedure must therefore terminate.
Its soundness does not require that the two tests detect every
possible obstruction to a realization.
\end{proof}
